\chapter{Einleitung}

\section{Versuchsziel}

Um den geometrischen Aufbau von Kristallstrukturen und die chemische Zusammensetzung von Materialien zerstörungsfrei zu untersuchen, können Röntgenstrahlen genutzt werden.

Es wird zunächst mithilfe von Bragg-Reflexion an einem NaCl-Kristall die charakteristische Röntgenstrahlung einer Röntgenröhre mit unbekannten Anodenmaterial untersucht. Das Material der Anode wird so bestimmt. Außerdem wird mit einer Langzeitmessung der Bragg-Reflexion die Feinstruktur $K_\alpha$-Linie der charakteristischen Röntgenstrahlung einer Molybdän-Anode untersucht.
Danach wird die chemische Zusammensetzung von unbekannten Materialien mithilfe einer Röntgenfluoreszenz-Messung zerstörungsfrei bestimmt. Dazu wird ein Röntgenenergiedetektor aus PiN-Photodiode und nachgeschalteten Vielkanalanalysator zur Messung der charakteristischen Röntgenfluoreszenzstrahlung des zu untersuchenden Materials genutzt. Aus den Messungen werden die Bestandteile des Materials sowie ihre jeweiligen relativen Massenanteile bestimmt.
Zuletzt wird mithilfe des Laue-Verfahrens die Gitterstruktur eines NaCl-Kristalls untersucht. Hierbei wird das Beugungsmuster der an den verschiedenen Gitterebenen des Kristalls reflektierten Röntgenstrahlung auf einem Röntgenfilm aufgenommen, der anschließend entwickelt wird. Aus dem Beugungsmuster werden schließlich die Abstände der reflektierenden Gitterebenen sowie die Wellenlänge der Röntgenstrahlung bestimmt \cite[1-2]{Versuchsanleitung428}.


\section{Theoretische Grundlagen der Röntgenstrahlung}

Das Energiespektrum von Röntgenstrahlung besteht aus einem kontinuierlichen Bremsstrahlungsteil und einem charakteristischen Spektrum aus diskreten Linien \cite[254]{Demtroeder3}.
Bremsstrahlung entsteht bei Geschwindigkeits- oder Richtungsänderungen von Elektronen mit kinetischer Energie im Bereich von keV bis MeV in Materie. Die dabei nach \uppercase{Hertz} entstehenden elektromagnetischen Wellen haben eine kontinuierliche Energieverteilung (bzw. Wellenlängenverteilung).

Um energiereiche Elektronen zu erzeugen, wird in der Regel eine Röntgenröhre genutzt, deren Aufbau in Abbildung \ref{fig:roehre_aufbau_demtr3_254} schematisch dargestellt ist.

\jafps{roehre_aufbau_demtr3_254}{Aufbau einer Röntgenröhre, entnommen aus \cite[254]{Demtroeder3}. An der Kathode K treten als Glühemission durch die angelegte Heizspannung $U_H$ Elektronen in die evakuierte Röhre aus, die anschließend durch die Beschleunigungsspannung $U$ zur Anode beschleunigt werden. Sie treffen in erster Näherung mit der Energie $E=eU$ auf die Anode auf, wo sie durch Wechselwirkung mit dem Anodenmaterial Röntgenstrahlung erzeugen, die durch ein Fenster aus der Röhre austritt.}{0.6}

Die maximale Energie, die Bremsstrahlung erreichen kann, ist die vollständige kinetische Energie der Elektronen, das Spektrum fällt also bei $E=eU$ auf eine Intensität von Null ab.

Außerdem weist das Spektrum von Röntgenstrahlung charakteristische Strahlung auf. Hierbei werden durch die einfallenden Elektronen in den Molekülen oder schweren Atomen des Anodenmaterials Hüllenelektronen aus tiefen (stark gebundenen) Zuständen ionisiert oder in freie höhere Energiezustände angeregt. In die so frei gewordenen niedrigen Elektronenzustände können Elektronen aus weniger stark gebundenen, höheren Zuständen übergehen, die Energiedifferenz $\Delta E = E_f - E_i $ der beiden Zustände wird als Photon mit Frequenz $\nu = \frac{\Delta E}{h}$ aus dem Molekül/Atom emittiert.
Die möglichen Zustandsübergänge und damit die diskreten emittierten Röntgenstrahlungsfrequenzen sind charakteristisch für die beteiligten Moleküle/Atome, aus ihrem Auftreten in einem Spektrum kann daher auf die Zusammensetzung des emittierenden Materials geschlossen werden \cite[253-256]{Demtroeder3}.
